\documentclass{article}
\usepackage{iclr2027_conference,times}
\usepackage[T1]{fontenc}
\usepackage{lmodern}
\usepackage{amsmath,amssymb,amsthm,mathtools}
\usepackage{booktabs}
\usepackage{microtype}
\usepackage{xcolor}
\usepackage{hyperref}
\usepackage{url}

\newcommand{\R}{\mathbb{R}}
\newcommand{\Sph}{\mathbb{S}}
\newcommand{\ip}[2]{\langle #1,#2\rangle}
\newcommand{\OPT}{\operatorname{OPT}}
\newcommand{\rad}{\operatorname{rad}}
\newcommand{\diam}{\operatorname{diam}}
\newcommand{\conv}{\operatorname{conv}}
\newcommand{\width}{\operatorname{width}}
\newcommand{\cE}{\mathcal{E}}
\newcommand{\cV}{\mathcal{V}}
\newcommand{\cF}{\mathcal{F}}
\newcommand{\cQ}{\mathcal{Q}}
\newcommand{\na}{\mathrm{na}}

\newtheorem{theorem}{Theorem}
\newtheorem{proposition}[theorem]{Proposition}
\newtheorem{lemma}[theorem]{Lemma}
\newtheorem{corollary}[theorem]{Corollary}
\theoremstyle{definition}
\newtheorem{definition}[theorem]{Definition}
\newtheorem{remark}[theorem]{Remark}

\title{The Price of Obliviousness in Noisy Linear Reconstruction}

\author{\textbf{Pahan Dewasurendra} \\ Johns Hopkins University}

\hypersetup{pdftitle={The Price of Obliviousness in Noisy Linear Reconstruction},pdfauthor={Pahan Dewasurendra},pdfkeywords={price, obliviousness, noisy, linear, reconstruction, unknown, point, queried, through},pdfsubject={preprint}}
\begin{document}

\maketitle

\begin{abstract}
An unknown point $x^\star\in\R^d$ is queried through unit linear functionals,
each answered with adversarial additive error at most $\delta$.  Recent work
determined the adaptive minimax error.  Its excess above the infinite-query
Jung limit $J_d\delta=\sqrt{2d/(d+1)}\,\delta$ vanishes doubly exponentially in
the number of queries, and it left the nonadaptive game open.  We characterize
that game.  For every fixed $d\geq2$, its minimax excess is
$\Theta_d(\delta T^{-2/(d-1)})$.  Thus adaptivity changes the convergence law
from polynomial to doubly exponential.  We also resolve the previously open
exponential-budget regime in high dimension.  If
$\log T/d\to\alpha\in(0,\infty)$, then the nonadaptive minimax error satisfies
\[
 \frac{\OPT^{\na}_d(T,\delta)}{\delta}
 \longrightarrow \sqrt{\frac{2}{1-e^{-2\alpha}}}.
\]
The upper bounds combine spherical coverings with Jung's theorem.  The lower
bounds use a common transcript that hides every vertex of a suitably rotated
and expanded regular simplex.  This construction is what makes classical
spherical-cap geometry sharp for point reconstruction.  Finally, a batched
strategy shows that each additional adaptive batch can double the polynomial
exponent, interpolating between the two convergence laws.
\end{abstract}

\section{Introduction}

Consider the following adversarial sensing problem.  A reconstructor seeks an
unknown $x^\star\in\R^d$.  On query $v\in\Sph^{d-1}$, it receives $r\in\R$
such that
\begin{equation}
    |r-\ip{v}{x^\star}|\leq\delta.                 \label{eq:model}
\end{equation}
After $T$ queries it outputs $\widehat x$ and incurs Euclidean error
$\|\widehat x-x^\star\|_2$.  The noise is bounded but otherwise adversarial.
This model isolates a basic question in adaptive sensing, robust localization,
and reconstruction from privacy-perturbed linear statistics
\citep{haupt2009adaptive,ariascatro2013fundamental,dinur2003revealing}.

\citet{filmus2026optimal} recently characterized the adaptive game.  With
$J_d=\sqrt{2d/(d+1)}$, its infinite-query error is $J_d\delta$.  For fixed
$d\geq2$, the excess above this limit is
$2^{-2^{\Theta_d(T)}}\delta$.  The mechanism has two phases.  An oblivious
spherical net first bounds the feasible diameter.  Later queries follow the
edges of a near-extremal simplex identified from previous answers.  Their paper
ends by asking what happens when every direction must be chosen before any
answer is observed.

This restriction is not cosmetic.  Nonadaptive queries are the natural model
when sensing directions must be deployed in parallel, when a fixed query batch
is audited before release, or when interaction is unavailable.  Classical
information-based complexity compares adaptive and nonadaptive recovery at a
coarse multiplicative scale.  For convex symmetric noisy information graphs,
the two errors differ by at most a factor of two
\citep{osipenko2019generalized}.  That comparison cannot see their convergence
to the same nonzero limit.  The open question is therefore an excess-risk
question.

\paragraph{Our results.}
We determine the nonadaptive excess in fixed dimension and its exact
high-dimensional phase curve.

\begin{theorem}[Fixed-dimensional nonadaptive rate]                 \label{thm:fixed}
For every fixed $d\geq2$, there are constants $0<c_d<C_d<\infty$ and
$T_d<\infty$ such that, for all $T\geq T_d$ and $\delta>0$,
\[
 c_d\delta T^{-2/(d-1)}
 \leq \OPT^{\na}_d(T,\delta)-J_d\delta
 \leq C_d\delta T^{-2/(d-1)}.
\]
\end{theorem}

Together with the adaptive result of \citet{filmus2026optimal}, Theorem
\ref{thm:fixed} gives a polynomial versus doubly-exponential separation in
excess error.  Equivalently, reaching excess at most $\epsilon\delta$ requires
$\Theta_d(\epsilon^{-(d-1)/2})$ nonadaptive queries, while adaptive query
complexity is $\Theta_d(\log\log(1/\epsilon))$.

\begin{table}[t]
\caption{Fixed-dimensional convergence.  Proper reconstructors output one
point.  Improper reconstructors predict each future linear functional and have
a different limiting benchmark.}
\label{tab:comparison}
\centering
\small
\begin{tabular}{llll}
\toprule
Setting & Output & Limiting error & Excess after $T$ queries \\
\midrule
Proper, adaptive & point & $J_d\delta$ & $2^{-2^{\Theta_d(T)}}\delta$ \\
Proper, nonadaptive & point & $J_d\delta$ & $\Theta_d(T^{-2/(d-1)})\delta$ \\
Improper, adaptive & function & $\delta$ & $\Theta_d(T^{-2/(d-1)})\delta$ \\
\bottomrule
\end{tabular}
\end{table}

The improper polynomial rate in the last row is due to
\citet{filmus2026optimal}.  It approaches the smaller benchmark $\delta$ with
an output that need not correspond to any point.  Its missed-direction lower
bound cannot prove a proper error near $J_d\delta$.  The simplex transcript in
this paper is needed to obtain the middle row.

The dimension dependence contains a second sharp phenomenon.  The existing
adaptive analysis identifies subexponential and superexponential regimes but
does not resolve budgets $T=\exp(\Theta(d))$.  For the nonadaptive game we get
an exact answer throughout this window.

\begin{theorem}[Exponential-budget phase curve]                     \label{thm:phase}
Let $T=T(d)$ satisfy $\log T/d\to\alpha\in(0,\infty)$.  Then for every
$\delta>0$,
\[
 \lim_{d\to\infty}\frac{\OPT^{\na}_d(T,\delta)}{\delta}
 = \sqrt{\frac{2}{1-e^{-2\alpha}}}.
\]
\end{theorem}

The curve diverges like $1/\sqrt{\alpha}$ as $\alpha\downarrow0$ and tends to
$\sqrt2$ only as $\alpha\to\infty$.  Hence a superexponential budget is
necessary and sufficient for vanishing nonadaptive excess as $d$ grows.

We also quantify how a small number of adaptive rounds improves the fixed-
dimensional exponent.  A $B$-batch strategy chooses all queries in a batch
before seeing any answer from that batch, while later batches can use earlier
answers.

\begin{theorem}[Batch upper hierarchy]                              \label{thm:batch}
Fix $d\geq2$ and $B\geq1$, and put $q_d=\binom{d+1}{2}$.  There are constants
$C_{d,B},M_{d,B}<\infty$ such that, whenever
$M=T-(B-1)q_d\geq M_{d,B}$,
\[
 \OPT^{(B)}_d(T,\delta)-J_d\delta
 \leq C_{d,B}\delta M^{-2^B/(d-1)}.
\]
\end{theorem}

The first batch is a spherical net.  Each later batch queries all edges of one
simplex, using the robust Jung refinement of \citet{filmus2026optimal}.  We do
not claim that Theorem \ref{thm:batch} is minimax optimal for intermediate $B$.

\paragraph{Proof idea.}
For fixed nonadaptive directions $V=\{v_1,\ldots,v_T\}$, define
\[
 \kappa(V)=\min_{u\in\Sph^{d-1}}\max_{i\leq T}|\ip{u}{v_i}|.
\]
This is the inradius of $\conv(\{\pm v_i\}_{i=1}^T)$.  Every transcript has
feasible diameter at most $2\delta/\kappa(V)$.  Jung's theorem then gives error
at most $J_d\delta/\kappa(V)$.  Good spherical covers make $\kappa(V)$ large.

The converse needs more than a missed direction.  A two-point ambiguity only
gives a limiting error $\delta$, not $J_d\delta$.  Instead take a regular
$d$-simplex of side $2\delta$.  Rotate it so that none of its
$\binom{d+1}{2}$ edge directions is too close to any query, then expand it.
Every queried directional width remains at most $2\delta$.  Answer each query
with the midpoint of the extreme vertex projections.  All $d+1$ expanded
vertices now produce one transcript, so one of them is at least the expanded
circumradius from the estimate.  A cap union bound finds the rotation.  In
fixed dimension it yields an avoided angle of order $T^{-1/(d-1)}$.  With
$T=\exp(\alpha d+o(d))$, the spherical-cap large-deviation exponent yields the
exact threshold $\sqrt{1-e^{-2\alpha}}$.

\section{Model and geometric reduction}                            \label{sec:model}

All strategies in this paper are deterministic, matching
\citet{filmus2026optimal}.  A nonadaptive strategy chooses
$v_1,\ldots,v_T\in\Sph^{d-1}$ before the game.  The adversary knows the
strategy, chooses $x^\star$, and returns any answers obeying
Equation \eqref{eq:model}.  The strategy then maps the transcript to an estimate.
The minimax error is
\[
 \OPT^{\na}_d(T,\delta)
 =\inf_{V,A}\sup_{x^\star\in\R^d}
   \sup_{|r_i-\ip{v_i}{x^\star}|\leq\delta}
       \|A(r_1,\ldots,r_T)-x^\star\|_2.
\]
The case $d=1$ is degenerate.  One query attains error $\delta$, so we assume
$d\geq2$.  Scaling gives
$\OPT^{\na}_d(T,\delta)=\delta\OPT^{\na}_d(T,1)$.

For a transcript $r$, let
\begin{equation}
 \Phi(V,r)=\{x\in\R^d:|\ip{v_i}{x}-r_i|\leq\delta
                    \text{ for all }i\leq T\}.                    \label{eq:feasible}
\end{equation}
For a bounded set $K$, let $\rad(K)=\inf_z\sup_{x\in K}\|x-z\|_2$.
The optimal output for a fixed transcript is a Chebyshev center of its feasible
set.  Thus the error of a design is $\sup_r\rad(\Phi(V,r))$ over nonempty
feasible sets.

\begin{proposition}[Covering upper bound]                           \label{prop:upper}
For every $V$ with $\kappa(V)>0$,
\[
 \sup_{r:\Phi(V,r)\neq\varnothing}\rad(\Phi(V,r))
 \leq \frac{J_d\delta}{\kappa(V)}.
\]
\end{proposition}

\begin{proof}
Take $x,y\in\Phi(V,r)$ and put $u=(x-y)/\|x-y\|_2$.  Some $v_i$
satisfies $|\ip{u}{v_i}|\geq\kappa(V)$.  Feasibility gives
$|\ip{x-y}{v_i}|\leq2\delta$, hence
$\|x-y\|_2\leq2\delta/\kappa(V)$.  Jung's theorem
\citep{jung1901kleinste,gruber2007convex} states
\[
 \rad(K)\leq\sqrt{\frac{d}{2(d+1)}}\diam(K).
\]
Substitution proves the claim.
\end{proof}

We next isolate the ambiguity construction used in both lower bounds.  For a
finite set $K$, write
$\width_K(v)=\max_{x\in K}\ip{v}{x}-\min_{x\in K}\ip{v}{x}$.

\begin{lemma}[Rotated-simplex transcript]                           \label{lem:simplex}
Let $V\subset\Sph^{d-1}$ and let $\Delta$ be a regular $d$-simplex of side
$2\delta$.  Suppose a rotation $R$ and $m\in(0,1]$ satisfy
\[
 |\ip{v}{Re}|\leq m
 \quad\text{for every }v\in V
 \text{ and every unit edge direction }e\text{ of }\Delta.
\]
Then the worst-case error of the design $V$ is at least $J_d\delta/m$.
\end{lemma}

\begin{proof}
Let $K=m^{-1}R\Delta$.  A directional width of a polytope is attained by a
pair of vertices.  Since every edge of $\Delta$ has length $2\delta$, the
hypothesis gives $\width_K(v)\leq2\delta$ for every $v\in V$.  Return
\[
 r_v=\tfrac12\big(\max_{x\in K}\ip{v}{x}
                  +\min_{x\in K}\ip{v}{x}\big).
\]
Every vertex of $K$ is consistent with every answer.  The Chebyshev radius of
these vertices is the circumradius $J_d\delta/m$.  Therefore some vertex is at
least that far from the design's output.  Because $V$ is fixed before the game,
the adversary can choose this vertex as $x^\star$ at the outset and return the
displayed common transcript.
\end{proof}

Lemma \ref{lem:simplex} captures why noncentral information fibers matter.
The zero-answer fiber is centrally symmetric and yields only a two-point lower
bound.  The least favorable fiber here contains a regular simplex.  This is
also why general factor-two adaptivity bounds in optimal recovery do not settle
the excess rate.

\section{Fixed dimension}                                          \label{sec:fixed}

For unit vectors, let
$\rho(u,v)=\arccos|\ip{u}{v}|$ be projective angular distance.  A projective
$\eta$-net has a point within distance $\eta$ of every line.  Standard
volumetric covering estimates give the following facts for fixed $d\geq2$:
\begin{equation}
 \begin{split}
 &\text{there is a projective $\eta$-net of size at most }
       a_d\eta^{-(d-1)},\\
 &\mu_d\{u:\rho(u,e)\leq\eta\}
       \leq b_d\eta^{d-1}\quad\text{for }0<\eta\leq\eta_d.
                                                               \label{eq:caps-fixed}
 \end{split}
\end{equation}
where $\mu_d$ is uniform spherical measure.  Appendix
\ref{app:caps} gives a short proof with explicit integral bounds.

\begin{proof}[Proof of Theorem \ref{thm:fixed}]
For the upper bound, choose $T$ directions forming a projective net of radius
$\eta\leq A_dT^{-1/(d-1)}$.  Then
$\kappa(V)\geq\cos\eta$.  Proposition \ref{prop:upper} and
$1/\cos\eta=1+\Theta(\eta^2)$ give
\[
 \OPT^{\na}_d(T,\delta)
 \leq \frac{J_d\delta}{\cos\eta}
 \leq J_d\delta+C_d\delta T^{-2/(d-1)}.
\]

For the lower bound, fix arbitrary queries $V$.  Independently rotate a regular
simplex $\Delta$.  Each of its $N_d=\binom{d+1}{2}$ unoriented edge directions
is marginally uniform on projective space.  A union bound and
Equation \eqref{eq:caps-fixed} show
\[
 \Pr\big\{\rho(Re,v)<\eta\text{ for some edge }e,v\in V\big\}
 \leq N_dTb_d\eta^{d-1}.
\]
Choose $\eta=c'_dT^{-1/(d-1)}$ so that the right side is below one.
Some rotation then has every edge-query angle at least $\eta$.  Lemma
\ref{lem:simplex} with $m=\cos\eta$ gives
\[
 \OPT^{\na}_d(T,\delta)
 \geq\frac{J_d\delta}{\cos\eta}
 \geq J_d\delta+c_d\delta T^{-2/(d-1)}.
\]
The Taylor inequalities used above hold after increasing $T_d$ if needed.
\end{proof}

The exponent has a simple geometric origin.  Covering a $(d-1)$-dimensional
projective sphere gives angular resolution $T^{-1/(d-1)}$.  Both the secant
loss in the upper bound and the safe expansion of the hidden simplex are
quadratic in that angle.

\section{The exponential-budget phase curve}                       \label{sec:phase}

The high-dimensional result uses a classical spherical-cap exponent.  If
$U$ is uniform on $\Sph^{d-1}$ and $m\in(0,1)$ is fixed, define
\[
 p_d(m)=\Pr\{\ip{U}{e_1}\geq m\},\qquad
 I(m)=-\tfrac12\log(1-m^2).
\]
Since $U_1^2$ has a $\operatorname{Beta}(1/2,(d-1)/2)$ distribution,
\begin{equation}
 \lim_{d\to\infty}\frac1d\log p_d(m)
 =\tfrac12\log(1-m^2)=-I(m).                    \label{eq:cap-ldp}
\end{equation}
The upper bound also needs the matching covering exponent.

\begin{lemma}[Spherical covering exponent]                          \label{lem:cover-exp}
For fixed $m\in(0,1)$, $\Sph^{d-1}$ can be covered by
$\exp(dI(m)+o(d))$ one-sided caps
$\{u:\ip{u}{v}\geq m\}$.
\end{lemma}

\begin{proof}
The cap has normalized area $p_d(m)$.  The spherical covering theorem of
\citet{boroczky2003covering}, sharpened by \citet{dumer2007covering}, gives a
cover of density at most $O(d\log d)$ for every acute cap.  Thus the number of
caps is at most $O(d\log d)/p_d(m)$.  Equation \eqref{eq:cap-ldp} proves the
claim.  This exponent is also the classical asymptotic inradius law for
few-vertex polytopes \citep{barany1988approximation,barany2005note}.
\end{proof}

\begin{proof}[Proof of Theorem \ref{thm:phase}]
Let
\[
 m_\alpha=\sqrt{1-e^{-2\alpha}},
 \qquad I(m_\alpha)=\alpha.
\]
For any fixed $m<m_\alpha$, Lemma \ref{lem:cover-exp} supplies at most $T$
query directions for all sufficiently large $d$.  These one-sided caps cover
the sphere, so $\kappa(V)\geq m$.  Proposition \ref{prop:upper} yields
\[
 \limsup_{d\to\infty}\frac{\OPT^{\na}_d(T,\delta)}{\delta}
 \leq\frac{\sqrt2}{m}.
\]
Letting $m\uparrow m_\alpha$ gives the desired upper limit.

For the converse, fix arbitrary $T$ queries and choose a uniform random
rotation of a regular simplex.  For any fixed $m>m_\alpha$, a union bound over
$N_d=\binom{d+1}{2}$ edges and $T$ queries gives
\[
 \Pr\big\{|\ip{Re}{v_i}|>m\text{ for some }e,i\big\}
 \leq 2N_dT p_d(m)
 =\exp\{d(\alpha-I(m))+o(d)\}\longrightarrow0.
\]
Hence a rotation with all correlations at most $m$ exists.  Lemma
\ref{lem:simplex} gives
\[
 \liminf_{d\to\infty}\frac{\OPT^{\na}_d(T,\delta)}{\delta}
 \geq\frac{\sqrt2}{m}.
\]
Letting $m\downarrow m_\alpha$ matches the upper limit.
\end{proof}

The argument uses classical optimal spherical coverage on the upper side and a
different object on the lower side.  The lower witness is not the uncovered
line itself.  It is one rotation of all simplex edges that simultaneously avoids
the queries.  The polynomial number $N_d$ of edges disappears in the exponential
scale, which is why the two exponents match exactly.

\section{A hierarchy from adaptive batches}                         \label{sec:batch}

The fixed-dimensional proof also clarifies how interaction enters.  We use one
result from the adaptive analysis of \citet{filmus2026optimal}.  In the
normalization $\delta=1/2$, if a feasible region has radius at least $J_d/2$ and
diameter at most $1+\beta$, their robust Jung theorem localizes its radius-
witnessing core to a $C_d\beta$ neighborhood of a regular unit simplex.  Querying
all $q_d$ edge directions of that simplex reduces the neighborhood radius from
$\gamma$ to at most $A_d\gamma^2$.  All $q_d$ directions depend only on the
previous feasible region, so they form one batch.

\begin{proof}[Proof of Theorem \ref{thm:batch}]
Use $M$ directions in the first batch to form a projective net of radius
$O_d(M^{-1/(d-1)})$.  Proposition \ref{prop:upper}, before applying Jung's
radius factor, gives feasible diameter
\[
 2\delta\big(1+O_d(M^{-2/(d-1)})\big).
\]
After scaling to $\delta=1/2$, robust Jung localization gives a simplex
neighborhood of radius
$\gamma_0=O_d(M^{-2/(d-1)})$.  For $b=1,\ldots,B-1$, spend one batch of
$q_d$ queries on the current simplex edges.  The refinement recurrence is
$\gamma_b\leq A_d\gamma_{b-1}^2$.  For fixed $d,B$ and sufficiently large
$M$, iteration yields
\[
 \gamma_{B-1}\leq C_{d,B}M^{-2^B/(d-1)}.
\]
The feasible radius is at most the regular-simplex radius plus
$\gamma_{B-1}$ in normalized units.  Rescaling proves the theorem.
\end{proof}

For fixed $B$, Theorem \ref{thm:batch} implies query complexity
$O_{d,B}(\epsilon^{-(d-1)/2^B})+(B-1)q_d$ for excess
$\epsilon\delta$.  The case $B=1$ matches Theorem \ref{thm:fixed}.  Once the
number of batches grows, the repeated squaring transitions to the adaptive
$\log\log(1/\epsilon)$ law.  A matching minimax characterization for every
intermediate batch budget remains open.

\section{Related work and limitations}                              \label{sec:related}

\paragraph{Linear reconstruction and adaptive sensing.}
Our direct predecessor is \citet{filmus2026optimal}, which introduced this
continuous linear-query game and proved the adaptive results summarized above.
It builds on approximate distance reconstruction
\citep{moran2025reconstruction}.  Adaptive sensing more commonly assumes a
structured signal and stochastic noise
\citep{haupt2009adaptive,ariascatro2013fundamental}.  Those goals and noise
models differ from worst-case localization of an unrestricted point.

\paragraph{Optimal recovery and bounded-error estimation.}
Optimal recovery studies algorithms from exact or inaccurate information
\citep{traub1988information,osipenko2019generalized}.  Recent sharp results
construct optimal linear recovery maps for prescribed operators under
hyperellipsoidal model classes and $\ell_1$ or $\ell_2$ data uncertainty
\citep{foucart2024radius}.  Set-membership and Chebyshev-center estimators
likewise recover from intersections of deterministic constraint sets
\citep{eldar2008minimax}.  This literature motivates our feasible-radius
formulation.  It typically fixes the information operator or imposes a bounded
model class.  The present question optimizes unit directions, compares answer-
dependent and oblivious designs, and asks for the excess above a nonzero Jung
limit.

\paragraph{Spherical covering.}
The projective-net estimates behind our upper bounds are classical.  Sharp
covering densities and equivalent inradius estimates appear in
\citet{boroczky2003covering,dumer2007covering,barany1988approximation}.  We use
their covering exponent rather than claim a new polytope approximation result.
The reconstruction lower bound is different: it must place an expanded simplex
inside one information fiber, which requires simultaneous avoidance by every
edge direction.

\paragraph{Limitations.}
We follow the predecessor's deterministic minimax convention.  The results also
cover public randomization when the realized design is visible before the secret
is selected.  Private randomized nonadaptive designs under alternative move
orders require a separate formulation.  Our estimator is information-theoretic
because it outputs an exact Chebyshev center.  Computing that center is generally
nontrivial, although standard convex optimization gives approximations.  Finally,
Theorem \ref{thm:batch} is an upper interpolation, not a matching batch lower
bound.

\section{Conclusion}

Nonadaptive noisy linear reconstruction has the same infinite-query Jung limit
as adaptive reconstruction, but it approaches that limit at a fundamentally
different speed.  A spherical net gives the optimal polynomial exponent, while
a rotated regular simplex supplies the matching ambiguity.  In high dimension,
the same construction converts the classical cap exponent into an exact
exponential-budget minimax curve.  These results resolve the nonadaptive problem
posed by \citet{filmus2026optimal} and isolate the geometric value of interaction:
adaptive batches repeatedly square the residual simplex uncertainty.

\section*{Reproducibility statement}
All results are proved analytically.  The supplement includes a deterministic
diagnostic that constructs regular simplices, searches rotations, verifies the
common-transcript inequalities, and checks finite-dimensional convergence of the
spherical-cap exponent.  These computations are sanity checks and are not used
as proof steps.

\bibliography{references}
\bibliographystyle{iclr2027_conference}

\appendix
\section{Spherical caps and projective nets}                        \label{app:caps}

This appendix records the cap estimates used in Sections
\ref{sec:fixed} and \ref{sec:phase}.  They are included to make the dependence
on the two asymptotic regimes explicit.

\subsection{Fixed dimension}

Let $C_d(\eta)=\{u\in\Sph^{d-1}:\arccos|\ip{u}{e_1}|\leq\eta\}$ be a
projective cap.  For $0<\eta<\pi/2$, normalized surface measure gives
\begin{equation}
 \mu_d(C_d(\eta))
 =\frac{2\int_0^\eta\sin^{d-2}\theta\,d\theta}
        {\int_0^\pi\sin^{d-2}\theta\,d\theta}.                    \label{eq:projective-cap}
\end{equation}
For fixed $d$, $\sin\theta\leq\theta$ yields
\[
 \mu_d(C_d(\eta))\leq c_d\eta^{d-1}.
\]
The reverse inequality with another constant holds for
$0<\eta\leq\eta_d$, since $\sin\theta\geq2\theta/\pi$ on
$[0,\pi/2]$.

Take a maximal set of projective points separated by more than $\eta$.
Maximality makes it an $\eta$-net.  Its projective caps of radius $\eta/2$
are disjoint.  The lower cap-volume estimate therefore bounds its cardinality
by $a_d\eta^{-(d-1)}$.  This proves both assertions in
Equation \eqref{eq:caps-fixed}.

\subsection{High-dimensional cap exponent}

For $U$ uniform on $\Sph^{d-1}$, the first coordinate has density
\begin{equation}
 f_d(t)=c_d(1-t^2)^{(d-3)/2},\qquad
 c_d=\frac{\Gamma(d/2)}{\sqrt\pi\,\Gamma((d-1)/2)}.                \label{eq:coord-density}
\end{equation}
Stirling's formula gives $\log c_d=o(d)$.  For fixed $m\in(0,1)$,
\[
 p_d(m)=\int_m^1 f_d(t)\,dt
 \leq c_d(1-m)(1-m^2)^{(d-3)/2}.
\]
This proves
\[
 \limsup_{d\to\infty}\frac1d\log p_d(m)
 \leq\tfrac12\log(1-m^2).
\]
For every fixed $\epsilon\in(0,1-m)$,
\[
 p_d(m)\geq
 \epsilon c_d\big(1-(m+\epsilon)^2\big)^{(d-3)/2}.
\]
First let $d\to\infty$, then let $\epsilon\downarrow0$.  This gives the
matching lower limit and proves Equation \eqref{eq:cap-ldp}.

For completeness, we spell out why the covering theorem supplies the same
exponent.  A cap $\{u:\ip{u}{v}\geq m\}$ has acute angular radius
$\phi=\arccos m$.  It can be represented as the intersection of a sphere of
radius $r=1/\sin\phi>1$ with a unit Euclidean ball.  Theorem 1 of
\citet{dumer2007covering}, with sphere dimension $d-1$, gives a cover whose
density is at most $O(d\log d)$, uniformly over $r>1$.  Dividing this density
by the normalized cap area $p_d(m)$ bounds the number of caps by
\[
 \frac{O(d\log d)}{p_d(m)}=\exp(dI(m)+o(d)).
\]

\section{Details for the nonadaptive lower bounds}                 \label{app:lower}

We verify the simplex normalizations and the move order used in Lemma
\ref{lem:simplex}.

\subsection{Regular-simplex normalization}

Let $z_0,\ldots,z_d\in\R^d$ obey
\[
 \ip{z_i}{z_i}=\frac{d}{d+1},\qquad
 \ip{z_i}{z_j}=-\frac1{d+1}\quad(i\neq j).
\]
These are centered regular-simplex vertices with side length $\sqrt2$ and
circumradius $\sqrt{d/(d+1)}$.  Multiplying each vertex by $\sqrt2\delta$
gives side $2\delta$ and circumradius
\[
 \sqrt2\delta\sqrt{\frac d{d+1}}=J_d\delta.
\]
After rotation and scaling by $1/m$, the circumradius is $J_d\delta/m$.

For any finite vertex set $K$ and direction $v$,
\[
 \width_K(v)=\max_{x,y\in K}\ip{v}{x-y}.
\]
Every difference of two regular-simplex vertices is an edge vector.  Thus the
edge-correlation condition in Lemma \ref{lem:simplex} implies
$\width_{m^{-1}R\Delta}(v)\leq2\delta$.

\subsection{One transcript and one secret}

Fix a deterministic nonadaptive algorithm.  Its complete query set $V$ is a
known function of the algorithm, independent of the secret and answers.  The
rotation in Lemma \ref{lem:simplex}, the scaled simplex $K$, the midpoint
answers, and the algorithm's output on those answers can therefore all be
determined before the interaction begins.  Choose a farthest vertex
\[
 x^\star\in\arg\max_{x\in K}\|x-A(r)\|_2.
\]
The midpoint construction ensures
$|r_v-\ip{v}{x^\star}|\leq\delta$ for every query.  Hence this is a legal
choice of one secret at the start of the game, not a post hoc switch between
secrets.  Since the minimum enclosing ball of a regular simplex is its
circumsphere,
\[
 \|x^\star-A(r)\|_2\geq\rad(K)=\frac{J_d\delta}{m}.
\]

\subsection{Why random rotation is enough}

Let $E$ be the set of $N_d=\binom{d+1}{2}$ unoriented unit edge directions of
a reference simplex.  If $R$ is Haar uniform on the orthogonal group, then for
each fixed $e\in E$, $Re$ is uniform on the sphere.  The different rotated
edges are dependent, but both lower bounds use only the union bound.  In fixed
dimension,
\[
 \Pr\{\rho(Re,v_i)<\eta\text{ for some }e,i\}
 \leq N_dT\mu_d(C_d(\eta)).
\]
In the exponential-dimensional regime,
\[
 \Pr\{|\ip{Re}{v_i}|>m\text{ for some }e,i\}
 \leq2N_dTp_d(m).
\]
Whenever the displayed upper bound is below one, at least one deterministic
rotation has the required simultaneous avoidance property.

\section{Batch refinement details}                                 \label{app:batch}

We give a more explicit derivation of Theorem \ref{thm:batch} from the robust
Jung machinery of \citet{filmus2026optimal}.  Their proof uses noise
$\delta=1/2$, so queried slabs have width one.  Write
\[
 \eta_M=a_dM^{-1/(d-1)},\qquad
 \beta_M=\frac1{\cos\eta_M}-1=O_d(M^{-2/(d-1)}).
\]
After querying the first-batch net, every feasible region $\Phi_0$ has
$\diam(\Phi_0)\leq1+\beta_M$.

The robust Jung theorem in the witness-core form states that, whenever
$\beta_M$ is below a dimension-dependent threshold and the feasible radius has
not already dropped below the Jung benchmark, there is a regular unit simplex
$\Delta_0$ such that
\[
 \operatorname{Witness}(\Phi_0)
 \subseteq\Delta_0+B(C_d\beta_M).
\]
The witness core has the same Chebyshev radius as the feasible region.  Put
$\gamma_0=C_d\beta_M$.

Suppose at the start of a later batch that
\[
 \operatorname{Witness}(\Phi_b)
 \subseteq\Delta_b+B(\gamma_b),
\]
where $\Delta_b$ is a regular unit simplex.  Query all its
$q_d=\binom{d+1}{2}$ unit edge directions in parallel.  The refinement argument
in the proof of Main Theorem 2 of \citet{filmus2026optimal} gives
\[
 \diam(\operatorname{Witness}(\Phi_{b+1}))
 \leq1+4\gamma_b^2.
\]
A second application of robust Jung gives another regular unit simplex and
\[
 \gamma_{b+1}\leq4C_d\gamma_b^2.
\]
For fixed $d,B$, take $M$ large enough that every iterate remains below the
robust-Jung threshold.  Induction yields
\[
 \gamma_{B-1}
 \leq C_{d,B}\beta_M^{2^{B-1}}
 \leq C'_{d,B}M^{-2^B/(d-1)}.
\]
The radius of $\Delta_{B-1}+B(\gamma_{B-1})$ is at most its simplex
circumradius plus $\gamma_{B-1}$.  This proves the normalized result.  General
$\delta$ follows by scaling.

If the radius ever drops below the Jung benchmark during the argument, the
desired upper bound is already true after enlarging $C_{d,B}$.  Thus the stated
case distinction in the robust theorem causes no gap.

\section{Endpoint consequences of the phase curve}                 \label{app:endpoints}

Theorem \ref{thm:phase} is deliberately stated for a fixed
$\alpha\in(0,\infty)$, where all cap estimates are uniform after fixing a small
margin around $m_\alpha$.  Two endpoint consequences follow without a uniform
large-deviation claim.

First, if $\log T/d\to\infty$, fix any $m<1$.  Lemma
\ref{lem:cover-exp} shows that the sphere can eventually be covered using at
most $T$ caps of threshold $m$.  Proposition \ref{prop:upper} gives
\[
 \limsup_{d\to\infty}\frac{\OPT^{\na}_d(T,\delta)}{\delta}
 \leq\frac{\sqrt2}{m}.
\]
Let $m\uparrow1$.  The universal Jung lower limit from
\citet{filmus2026optimal} gives equality with $\sqrt2$.

Second, if $\log T=o(d)$, nonadaptive strategies are a subset of adaptive
strategies.  The subexponential lower bound of \citet{filmus2026optimal}
therefore implies
\[
 \OPT^{\na}_d(T,\delta)-J_d\delta\longrightarrow\infty.
\]
The limiting curve in Theorem \ref{thm:phase} is consistent with both
endpoints.

\section{Computational sanity checks}                               \label{app:checks}

The file \texttt{verify\_geometry.py} performs two independent diagnostics.

First, for $d\in\{2,3,5\}$ it constructs centered regular-simplex coordinates,
draws a fixed nonadaptive query set, searches over seeded random rotations, and
chooses the rotation with the smallest maximum edge-query correlation $m$.  It
then expands the simplex by $1/m$, forms every midpoint response, and checks
\[
 \max_i\width_K(v_i)\leq2\delta,
 \qquad
 \max_{x\in K,i}|\ip{v_i}{x}-r_i|\leq\delta,
 \qquad
 \rad(K)=\frac{J_d\delta}{m}.
\]
All inequalities pass to floating-point tolerance $10^{-10}$.

Second, it evaluates the exact two-sided cap tail using
$U_1^2\sim\operatorname{Beta}(1/2,(d-1)/2)$.  For several values of $\alpha$
and $d\leq500$, it solves for the finite-dimensional correlation whose cap
probability is $e^{-\alpha d}$ and compares it with
$\sqrt{1-e^{-2\alpha}}$.  The absolute discrepancy decreases monotonically in
every tested sequence.  The generated CSV files contain the raw values.

These checks are not invoked in any proof.  They guard the edge-length,
noise-width, and beta-tail normalizations that are easiest to misstate.


\end{document}
